% GENERATED FILE. Edit canonical lessons, then run npm run generate:books.
% canonical-source-hash: fnv1a64:687a15a805680fd2
% canonical-lessons: TE-C468-murugu, TE-C468-parisuddhyam, TE-C468-durvasana, TE-C468-murikivada, TE-C468-perukupovu

\chapter{Drains and Rubbish}
\label{ch:te-drains-and-rubbish}
\hlchaptermodality{468}

\begin{chapteropening}
\textbf{By the end of this chapter:} \emph{I can say sewage, sanitation, a stench, a slum and to pile up.}

Complete the last lesson of chapter 468: I can say sewage, sanitation, a stench, a slum and to pile up.
\end{chapteropening}

\section[\texorpdfstring{\te{మురుగు} (murugu)}{మురుగు (murugu)}]{\te{మురుగు} (murugu) — sewage; foul water}

\label{lesson:TE-C468-murugu}



\begin{quote}
Before the new one: say the Telugu for spacious, then the Telugu for deserted.
\end{quote}

\subsection*{You'll want to know: \te{మురుగు}}
\textbf{\te{మురుగు}} — \emph{murugu} — \textquotedblleft{}sewage; foul water\textquotedblright{}.

\begin{quote}
\te{వర్షం} \te{వల్ల} \te{మురుగు} \te{నీరు} \te{రోడ్డు} \te{మీదకు} \te{వచ్చింది}.
\end{quote}

\begin{quote}
\emph{varṣaṁ valla murugu nīru rōḍḍu mīdaku vaccindi.}
\end{quote}

\begin{quote}
Because of the rain sewage came up onto the road.
\end{quote}

\te{మురుగు} \te{నీరు}, sewage; \te{మురుగు} \te{కాలువ}, a drain. \te{మురికి}, dirt, is the same family.

\begin{cousinweb}
\textbf{\te{మురుగు}} (\emph{murugu}) is native Dravidian: DEDR 5007 lists Telugu \emph{muṟugu}, \textquotedblleft{}to become putrid\textquotedblright{}, beside \emph{muṟiki}, \textquotedblleft{}dirt, filth\textquotedblright{}.
\end{cousinweb}

\begin{grammarlens}[title={What you've built}]
One more word for getting around.
\end{grammarlens}

\subsection*{Your turn}
\begin{itemize}
\raggedright
  \item \emph{Say it:} \emph{murugu}
  \item \emph{Say it:} \emph{murugu}, once more
  \item \emph{Say it:} \emph{varṣaṁ valla murugu nīru rōḍḍu mīdaku vaccindi.} — Because of the rain sewage came up onto the road.
  \item \emph{Recall:} say the Telugu for spacious, then the Telugu for deserted, then say \emph{murugu} again
\end{itemize}

\begin{tcolorbox}[breakable,colback=teal!4,colframe=teal!35!black,title={Before you move on}]
What does \te{మురుగు} mean? (\textquotedblleft{}Sewage; foul water\textquotedblright{}.) Say it once more.
\end{tcolorbox}
\section[\texorpdfstring{\te{పారిశుద్ధ్యం} (pāriśuddhyaṁ)}{పారిశుద్ధ్యం (pāriśuddhyaṁ)}]{\te{పారిశుద్ధ్యం} (pāriśuddhyaṁ) — sanitation}

\label{lesson:TE-C468-parisuddhyam}



\begin{quote}
Before the new one: say the Telugu for deserted, then the Telugu for sewage.
\end{quote}

\subsection*{You'll want to know: \te{పారిశుద్ధ్యం}}
\textbf{\te{పారిశుద్ధ్యం}} — \emph{pāriśuddhyaṁ} — \textquotedblleft{}sanitation\textquotedblright{}.

\begin{quote}
\te{పురపాలక} \te{సంఘం} \te{పారిశుద్ధ్యం} \te{మీద} \te{శ్రద్ధ} \te{పెట్టాలి}.
\end{quote}

\begin{quote}
\emph{purapālaka saṅghaṁ pāriśuddhyaṁ mīda śraddha peṭṭāli.}
\end{quote}

\begin{quote}
The municipality must pay attention to sanitation.
\end{quote}

An official word for keeping a town clean: drains, rubbish, toilets.

\begin{cousinweb}
\textbf{\te{పారిశుద్ధ్యం}} (\emph{pāriśuddhyaṁ}) is built from Sanskrit \emph{pariśuddha}, \textquotedblleft{}cleaned, purified\textquotedblright{} (Monier-Williams), as an abstract noun: the state of being clean.
\end{cousinweb}

\begin{grammarlens}[title={What you've built}]
One more word for getting around.
\end{grammarlens}

\subsection*{Your turn}
\begin{itemize}
\raggedright
  \item \emph{Say it:} \emph{pāriśuddhyaṁ}
  \item \emph{Say it:} \emph{pāriśuddhyaṁ}, once more
  \item \emph{Say it:} \emph{purapālaka saṅghaṁ pāriśuddhyaṁ mīda śraddha peṭṭāli.} — The municipality must pay attention to sanitation.
  \item \emph{Recall:} say the Telugu for deserted, then the Telugu for sewage, then say \emph{pāriśuddhyaṁ} again
\end{itemize}

\begin{tcolorbox}[breakable,colback=teal!4,colframe=teal!35!black,title={Before you move on}]
What does \te{పారిశుద్ధ్యం} mean? (\textquotedblleft{}Sanitation\textquotedblright{}.) Say it once more.
\end{tcolorbox}
\section[\texorpdfstring{\te{దుర్వాసన} (durvāsana)}{దుర్వాసన (durvāsana)}]{\te{దుర్వాసన} (durvāsana) — a stench; a bad smell}

\label{lesson:TE-C468-durvasana}



\begin{quote}
Before the new one: say the Telugu for sewage, then the Telugu for sanitation.
\end{quote}

\subsection*{You'll want to know: \te{దుర్వాసన}}
\textbf{\te{దుర్వాసన}} — \emph{durvāsana} — \textquotedblleft{}a stench; a bad smell\textquotedblright{}.

\begin{quote}
\te{చెత్త} \te{కుప్ప} \te{నుంచి} \te{దుర్వాసన} \te{వస్తోంది}.
\end{quote}

\begin{quote}
\emph{cetta kuppa nuñci durvāsana vastōndi.}
\end{quote}

\begin{quote}
A stench is coming from the heap of rubbish.
\end{quote}

\te{దుర్}, bad, + \te{వాసన}, a smell. Its opposite is \te{సువాసన}, a fragrance.

\begin{cousinweb}
\textbf{\te{దుర్వాసన}} (\emph{durvāsana}) is built from Sanskrit \emph{dur-}, \textquotedblleft{}bad, ill\textquotedblright{}, and \emph{vāsana}, \textquotedblleft{}perfuming, scent\textquotedblright{} (Monier-Williams).
\end{cousinweb}

\begin{grammarlens}[title={What you've built}]
One more word for getting around.
\end{grammarlens}

\subsection*{Your turn}
\begin{itemize}
\raggedright
  \item \emph{Say it:} \emph{durvāsana}
  \item \emph{Say it:} \emph{durvāsana}, once more
  \item \emph{Say it:} \emph{cetta kuppa nuñci durvāsana vastōndi.} — A stench is coming from the heap of rubbish.
  \item \emph{Recall:} say the Telugu for sewage, then the Telugu for sanitation, then say \emph{durvāsana} again
\end{itemize}

\begin{tcolorbox}[breakable,colback=teal!4,colframe=teal!35!black,title={Before you move on}]
What does \te{దుర్వాసన} mean? (\textquotedblleft{}A stench; a bad smell\textquotedblright{}.) Say it once more.
\end{tcolorbox}
\section[\texorpdfstring{\te{మురికివాడ} (murikivāḍa)}{మురికివాడ (murikivāḍa)}]{\te{మురికివాడ} (murikivāḍa) — a slum}

\label{lesson:TE-C468-murikivada}



\begin{quote}
Before the new one: say the Telugu for sanitation, then the Telugu for a stench.
\end{quote}

\subsection*{You'll want to know: \te{మురికివాడ}}
\textbf{\te{మురికివాడ}} — \emph{murikivāḍa} — \textquotedblleft{}a slum\textquotedblright{}.

\begin{quote}
\te{నగరంలో} \te{చాలా} \te{మంది} \te{మురికివాడల్లో} \te{ఉంటారు}.
\end{quote}

\begin{quote}
\emph{nagaraṁlō cālā mandi murikivāḍallō uṇṭāru.}
\end{quote}

\begin{quote}
In the city a lot of people live in slums.
\end{quote}

Plural: \te{మురికివాడలు}, \emph{murikivāḍalu}.

\te{మురికి}, dirty, + \te{వాడ}, a quarter of a town: a crowded, poor part of a town. Plural \te{మురికివాడలు}.

\begin{cousinweb}
\textbf{\te{మురికివాడ}} (\emph{murikivāḍa}) is a Telugu compound: \te{మురికి}, \textquotedblleft{}dirt, filth\textquotedblright{}, native Dravidian (DEDR 5007), and \te{వాడ}, \textquotedblleft{}a quarter of a town, a settlement\textquotedblright{}, usually traced to Sanskrit \emph{vāṭa}, \textquotedblleft{}an enclosure\textquotedblright{}, through Prakrit \emph{vāḍa} (CDIAL 11480).
\end{cousinweb}

\begin{grammarlens}[title={What you've built}]
One more word for getting around.
\end{grammarlens}

\subsection*{Your turn}
\begin{itemize}
\raggedright
  \item \emph{Say it:} \emph{murikivāḍa}
  \item \emph{Say it:} \emph{murikivāḍa}, once more
  \item \emph{Say it:} \emph{nagaraṁlō cālā mandi murikivāḍallō uṇṭāru.} — In the city a lot of people live in slums.
  \item \emph{Recall:} say the Telugu for sanitation, then the Telugu for a stench, then say \emph{murikivāḍa} again
\end{itemize}

\begin{tcolorbox}[breakable,colback=teal!4,colframe=teal!35!black,title={Before you move on}]
What does \te{మురికివాడ} mean? (\textquotedblleft{}A slum\textquotedblright{}.) Say it once more.
\end{tcolorbox}
\section[\texorpdfstring{\te{పేరుకుపోవు} (pērukupōvu)}{పేరుకుపోవు (pērukupōvu)}]{\te{పేరుకుపోవు} (pērukupōvu) — to pile up; to accumulate}

\label{lesson:TE-C468-perukupovu}



\begin{quote}
Before the new one: say the Telugu for a stench, then the Telugu for a slum.
\end{quote}

\subsection*{You'll want to know: \te{పేరుకుపోవు}}
\textbf{\te{పేరుకుపోవు}} — \emph{pērukupōvu} — \textquotedblleft{}to pile up; to accumulate\textquotedblright{}.

\begin{quote}
\te{వీధుల్లో} \te{చెత్త} \te{పేరుకుపోయింది}.
\end{quote}

\begin{quote}
\emph{vīdhullō cetta pērukupōyindi.}
\end{quote}

\begin{quote}
Rubbish has piled up in the streets.
\end{quote}

Past (it …): \te{పేరుకుపోయింది}, \emph{pērukupōyindi}.

Said of rubbish, dust, work or debts left to grow. Its subject is a thing, so the past shown is the third person: \te{పేరుకుపోయింది}, it piled up.

\begin{cousinweb}
\textbf{\te{పేరుకుపోవు}} (\emph{pērukupōvu}) is native Dravidian: \te{పేరు}-, from \emph{pēr(u)cu}, \textquotedblleft{}to pile up, to heap\textquotedblright{} (DEDR 4446), with -\te{కొని} and \te{పోవు}, \textquotedblleft{}to go\textquotedblright{} (DEDR 4572).
\end{cousinweb}

\begin{grammarlens}[title={What you've built}]
One more word for getting around.
\end{grammarlens}

\subsection*{Your turn}
\begin{itemize}
\raggedright
  \item \emph{Say it:} \emph{pērukupōvu}
  \item \emph{Say it:} \emph{pērukupōvu}, once more
  \item \emph{Say it:} \emph{vīdhullō cetta pērukupōyindi.} — Rubbish has piled up in the streets.
  \item \emph{Recall:} say the Telugu for a stench, then the Telugu for a slum, then say \emph{pērukupōvu} again
\end{itemize}

\begin{tcolorbox}[breakable,colback=teal!4,colframe=teal!35!black,title={Before you move on}]
What does \te{పేరుకుపోవు} mean? (\textquotedblleft{}To pile up; to accumulate\textquotedblright{}.) Say it once more.
\end{tcolorbox}
